\documentclass[a4paper,12pt]{article}
\usepackage[margin=1.5cm]{geometry}
\usepackage{amsmath, amssymb}
\usepackage{array}
\usepackage{alphalph}

\usepackage{amsmath, amssymb, amsfonts}
\usepackage{multicol}
\usepackage{multirow}

\usepackage{enumitem}
\usepackage{booktabs}
\usepackage{tabularx}
\usepackage{array}
\usepackage{framed}
\usepackage{xcolor}
\usepackage{tikz}
\usepackage{fancyhdr}

% Estilos para cajas de teoría y recordatorios
\definecolor{graybox}{gray}{0.92}
\newenvironment{boxsummary}{%
  \def\FrameCommand{\fcolorbox{black}{graybox}}%
  \MakeFramed{\FrameRestore}%
  \small\textbf{RECORDAR:}\par\smallskip
}{%
  \endMakeFramed
}


\begin{document}

\section*{Tema 2: Potencias y Raíces}

\subsection*{1. Completa la siguiente tabla}

\begin{center}
\renewcommand{\arraystretch}{1.5}
\begin{tabular}{|c|c|c|c|c|}
\hline
\textbf{Potencia} & \textbf{Base} & \textbf{Exponente} & \textbf{Valor} & \textbf{Se lee} \\ \hline
$(-1)^{7}$ & & & & \\ \hline
$(-2)^{5}$ & & & & \\ \hline
$(+3)^{3}$ & & & & \\ \hline
$(-4)^{3}$ & & & & \\ \hline
$8^{2}$ & & & & \\ \hline
$(+7)^{0}$ & & & & \\ \hline
$-2^{2}$ & & & & \\ \hline
\end{tabular}
\end{center}

\vspace{0.5cm}

\subsection*{2. ¿Qué signo tienen las potencias siguientes?}

\begin{enumerate}[label=\alph*)]
    \item $6^{3}$
    \item $(-3)^{12}$
    \item $3^{21}$
    \item $(-3)^{21}$
    \item $(-2)^{4}$
    \item $5^{32}$
    \item $(-3)^{5}$
    \item $4^{51}$
    \item $(-1)^{17}$
    \item $3^{3}$
    \item $(-2)^{3}$
\end{enumerate}

\vspace{0.5cm}

\subsection*{3. Expresa estas operaciones como una sola potencia.}

\begin{enumerate}[label=\alph*.]
    \item $(5^{8} \cdot 5^{4}) : (5^{2})^{5} =$
    \item $[(-2)^{6} \cdot (+2)^{3}] : [(+2)^{3}]^{2} =$
    \item $(-12)^{-7} \cdot [(-3)^{5} \cdot 4^{5}] =$
    \item $25^{3} : [(-15)^{5} : 3^{5}] =$
    \item $6^{3} : [(2^{7} : 2^{6}) \cdot 3]^{-2} =$
    \item $\left[\left(\dfrac{2}{5}\right)^{8} : \left(\dfrac{2}{5}\right)^{3}\right]^{-2} : \left(\dfrac{4}{25}\right)^{-1} =$
    \item $5 : (+1)^{2} - 3 : \left(\dfrac{1}{2} - \dfrac{1}{4}\right) =$
    \item $\dfrac{2}{3}\left(\dfrac{3}{4} - \dfrac{1}{2}\right)^{2} - \dfrac{1}{6}\left(\dfrac{5}{6} - \dfrac{1}{3}\right)^{2} =$
    \item $8^{4} : (2^{5} \cdot 4^{2}) =$
    \item $[(6^{2})^{2} \cdot 4^{4}] : (2^{3})^{4} =$
    \item $[4^{-4}] \cdot (2^{3})^{4} =$
    \item $2^{3} : 8^{-2} \cdot 2^{4} =$
    \item $81^{5} \cdot 27^{-2} =$
    \item $\dfrac{3^{7} \cdot 9^{3} \cdot 3^{-3}}{81 \cdot 9^{-5}}$
\end{enumerate}

\subsection*{4. Calcula sin usar la calculadora el valor de las siguientes potencias.}

\begin{enumerate}[label=\alph*)]
    \begin{minipage}[t]{0.3\textwidth}
        \item $2^{5} =$
        \item $(-2)^{5} =$
        \item $3^{4} =$
        \item $(-3)^{4} =$
        \item $1^{5} =$
        \item $(-1)^{5} =$
        \item $(-1)^{6} =$
        \item $(-1)^{37} =$
        \item $3^{0} =$
    \end{minipage}
    \begin{minipage}[t]{0.3\textwidth}
        \item[j)] $(-2)^{2} =$
        \item[k)] $(-5)^{0} =$
        \item[l)] $(-2)^{4} =$
        \item[m)] $-2^{4} =$
        \item[n)] $(-3)^{3} =$
        \item[o)] $-3^{3} =$
        \item[p)] $1^{34} =$
        \item[q)] $(-1)^{56} =$
        \item[r)] $(-1)^{57} =$
    \end{minipage}
    \begin{minipage}[t]{0.3\textwidth}
        \item[s)] $\left(\dfrac{1}{2}\right)^{3} =$
        \item[t)] $\left(\dfrac{1}{3}\right)^{2} =$
        \item[u)] $9^{2} =$
        \item[v)] $(-9)^{2} =$
        \item[w)] $\left(\dfrac{3}{2}\right)^{2} =$
        \item[x)] $9^{3} =$
        \item[y)] $(-9)^{3} =$
        \item[z)] $0{,}4^{2} =$
        \item[$\alpha$)] $60^{2} =$
    \end{minipage}
\end{enumerate}

\vspace{0.8cm}

\textbf{Consecuencias:}
\begin{itemize}
    \item $(\text{nº negativo})^{\text{par}} =$
    \item $(\text{nº negativo})^{\text{impar}} =$
    \item $1^{n} =$
    \item $(-1)^{\text{par}} =$
    \item $(-1)^{\text{impar}} =$
\end{itemize}

\vspace{0.5cm}

\subsection*{5. Simplifica, mediante las propiedades de las potencias, dejando el resultado como entero o fracción (salvo si es muy elevado en cuyo caso puede dejarse como potencia), no vale usar calculadora.}

\begin{enumerate}[label=\alph*)]
    \item $\left[\left(\dfrac{5}{2}\right)^{3}\right]^{-4} \cdot \left(\dfrac{4}{5}\right)^{-2} = \hfill \text{(Soluc: } 2^{7}/5^{10}\text{)}$
    \vspace{0.2cm}
    \item $\left(\dfrac{6}{5}\right)^{6} \cdot \left(-\dfrac{10}{3}\right)^{-4} = \hfill \text{(Soluc: } 3^{10} \cdot 2^{2} / 5^{10}\text{)}$
    \vspace{0.2cm}
    \item $\dfrac{2^{-3} \cdot (-2)^{4} \cdot (-4)^{-1}}{-2} = \hfill \text{(Soluc: } 1/4\text{)}$
    \vspace{0.2cm}
    \item $(-1)^{3} + (-1)^{2} + (-1) = \hfill \text{(Soluc: } -1\text{)}$
    \vspace{0.2cm}
    \item $2 \cdot (-1)^{3} - 4 \cdot (-1)^{2} + 2 \cdot (-1) = \hfill \text{(Soluc: } -8\text{)}$
    \vspace{0.2cm}
    \item $\dfrac{\left(\dfrac{1}{2}\right)^{3} \cdot \left(-\dfrac{1}{4}\right)^{2}}{2^{-1}} = \hfill \text{(Soluc: } 1/128\text{)}$
    \vspace{0.2cm}
    \item $2 \cdot (-2)^{4} + 3 \cdot (-2)^{3} - 4 \cdot (-2)^{2} - 3 \cdot (-2) = \hfill \text{(Soluc: } -2\text{)}$
\end{enumerate}

\end{document}